\documentclass[a4paper,10.5pt]{article}
\usepackage[T1]{fontenc}\usepackage[utf8]{inputenc}
\usepackage[italian]{babel}
\usepackage[top=2.2cm,bottom=2.2cm,left=2.35cm,right=2.35cm]{geometry}
\usepackage{hyperref}\usepackage{xcolor}\definecolor{accent}{HTML}{1F4E78}\hypersetup{colorlinks=true,urlcolor=accent}
\pagestyle{empty}\setlength{\parindent}{0pt}\setlength{\parskip}{7pt}
\begin{document}
\textbf{Prabhavit Gupta}\\
Venezia, Italia \textbar{} +39 351 610 7807 \textbar{} prabhavitg@gmail.com\\[10pt]
Gentile team di NewCo AI,

vi scrivo per candidarmi al ruolo di Freelance AI Automation Engineer per il vostro nuovo progetto dedicato agli agenti AI e all'automazione dei processi per le piccole e medie imprese. Lavoro come libero professionista con Partita IVA attiva e negli ultimi tre anni ho costruito pipeline di automazione e integrazioni AI per diverse aziende europee, un percorso che si adatta bene a quello che state cercando.

Per Urlink, a Barcellona, ho lavorato come project manager e sviluppatore assistito da AI: ho costruito la piattaforma e la sua logica con sviluppo assistito da intelligenza artificiale e ho implementato automazioni su richiesta diretta del cliente, gestendo io stesso i requisiti insieme al cliente. Questo tipo di collaborazione da freelance, a stretto contatto con un unico referente, è molto simile a come descrivete il vostro progetto.

Per FlowSpace, un marketplace di spazi di coworking con sede a Londra, ho costruito e gestisco un bot AI che, a partire da una brochure in qualsiasi formato (PDF, sito, foto o foglio di calcolo), genera automaticamente un annuncio strutturato con verifica umana nel processo, usando n8n e le API di Claude. Ho anche sviluppato agenti di supporto AI che rispondono alle richieste dei clienti a partire dai dati live della piattaforma e passano il caso a una persona quando la confidenza è bassa, oltre a integrazioni API (OfficeRnD e Nexudus) per le prenotazioni automatizzate.

Uso regolarmente Python per automazioni, scraping e integrazioni con API di modelli LLM, n8n per orchestrare i workflow, ho una conoscenza di base di FastAPI e di Docker per containerizzare le applicazioni, e lavoro abitualmente con GNU/Linux e Git. Non ho esperienza diretta con Selenium, ma ho automatizzato test e processi di scraping con Playwright.

Sono cittadino indiano con status di soggiornante di lungo periodo UE in Italia, autorizzato a lavorare in Italia senza necessità di sponsorizzazione, e la collaborazione sarebbe gestita tramite la mia Partita IVA attiva. Sono disponibile da subito.

Cordiali saluti,\\
Prabhavit Gupta
\end{document}
